\section*{How the Proof Works}

\begingroup
\setlength{\parskip}{0.38em}

For the global three-dimensional Navier--Stokes regularity problem, the fluid
obeys
\[
\partial_tu+(u\!\cdot\nabla)u+\nabla p=\nu\Delta u,
\qquad
\nabla\!\cdot u=0,
\qquad
u(\cdot,0)=u_0.
\]
A strain in the flow pulls a vortex tube along its axis. As the vortex
stretches, incompressibility forces it to contract across its width. Each
material parcel keeps its volume, so lengthening the tube narrows its
cross-section and the rotating core inside it. That strain also acts on the
vorticity
\(\omega=\nabla\times u\). Curling the equation gives
\[
\partial_t\omega+(u\!\cdot\nabla)\omega
=(\omega\!\cdot\nabla)u+\nu\Delta\omega.
\]
The left side follows the change in rotation with the moving fluid. When the
strain is aligned with the vortex,
\(\omega\cdot((\omega\!\cdot\nabla)u)>0\), so the stretching contribution raises
\(|\omega|\) while the core narrows. Since \(\omega\) is one spatial derivative
of \(u\), its amplification means that velocity changes more sharply across the
core. The gradients rise.

As the core narrows, the nonlinear term can sharpen its velocity profile more
intensely, but viscosity also acts more strongly on that smaller width. Scaling
tells us how their relative strengths change: it compresses the spatial pattern
of a solution while adjusting its velocity and clock so the rescaled fields can
still satisfy the equation. Given a solution \((u,p)\) and a factor
\(\lambda>1\), define
\[
u_\lambda(x,t)=\lambda u(\lambda x,\lambda^2t),
\qquad
p_\lambda(x,t)=\lambda^2p(\lambda x,\lambda^2t).
\]
This sends a feature of width \(r\), velocity scale \(U\), and duration \(\tau\)
to one of width \(r/\lambda\), velocity scale \(\lambda U\), and duration
\(\tau/\lambda^2\). Substitution gives
\[
\begin{aligned}
&\partial_tu_\lambda
 +(u_\lambda\!\cdot\nabla)u_\lambda
 +\nabla p_\lambda-\nu\Delta u_\lambda\\
&\qquad
=\lambda^3
 \bigl[\partial_tu+(u\!\cdot\nabla)u+\nabla p-\nu\Delta u\bigr]
 (\lambda x,\lambda^2t)
=0.
\end{aligned}
\]
Every term acquires the factor \(\lambda^3\). Viscosity strengthens as the core
narrows, but the nonlinearity strengthens in exact proportion. Compression
preserves their balance.

The velocity and gradients can rise without the kinetic energy rising with
them. The compressed core occupies a volume smaller by \(\lambda^{-3}\), while
the square of the velocity contributes \(\lambda^2\). In three dimensions,
\[
\|u_\lambda(\cdot,t)\|_2^2
=\lambda^{-1}\|u(\cdot,\lambda^2t)\|_2^2.
\]
A finite energy bound can survive while the core becomes narrower and its
gradients become steeper.

To find a norm that neither loses nor gains size just because the core was
compressed, use the homogeneous Sobolev scale:
\[
\|u_\lambda(\cdot,t)\|_{\dot H^s}
=\lambda^{s-\frac12}\|u(\cdot,\lambda^2t)\|_{\dot H^s}.
\]
At \(s=0\), this recovers the loss of kinetic energy under compression. Below
\(s=\tfrac12\), the norm becomes smaller; above it, the norm becomes larger. At
\(s=\tfrac12\), the factor is one. This is the critical height: narrowing the
core does not automatically change its measured size there.

The factor \(\lambda^2\) in time now matters. Suppose an evolving core repeats
this scale change, contracting by a fixed factor
\(0<\rho<1\) at each stage. Then
\[
\begin{gathered}
r_j=\rho^jr_0,
\qquad U_j=\rho^{-j}U_0,
\qquad \tau_j=\rho^{2j}\tau_0,\\[3pt]
\displaystyle
\sum_{j=0}^{\infty}\tau_j
=\tau_0\sum_{j=0}^{\infty}\rho^{2j}
=\frac{\tau_0}{1-\rho^2}<\infty.
\end{gathered}
\]
The widths collapse and the velocity scale diverges, yet the durations have a
finite sum. Infinitely many narrower, faster stages fit inside a finite
interval: this is the Zeno cascade. Scaling does not prove that a solution must
follow this pattern. It shows why finite energy and the scale balance between
nonlinearity and viscosity do not rule it out.

\endgroup

In order for a singularity to form, velocity has to blow up. In order for
that to happen, its acceleration would also have to blow up, which means its
jerk has to blow up, and so on, and so on:
\[
u,\qquad \partial_tu,\qquad \partial_t^2u,\qquad \ldots
\]
What you'd expect to see in this pattern, if a singularity were to form, as
you go up the time derivatives, you'd expect to see less spatial regularity,
not more.

The equation moves in the other direction. Its first three time derivatives
are
\[
\begin{aligned}
\partial_tu
={}&\nu\Delta u
 -(u\!\cdot\nabla)u
 -\nabla p,\\[3pt]
\partial_t^2u
={}&\nu\Delta\partial_tu
 -(\partial_tu\!\cdot\nabla)u
 -(u\!\cdot\nabla)\partial_tu
 -\nabla\partial_tp,\\[3pt]
\partial_t^3u
={}&\nu\Delta\partial_t^2u
 -(\partial_t^2u\!\cdot\nabla)u
 -2(\partial_tu\!\cdot\nabla)\partial_tu
 -(u\!\cdot\nabla)\partial_t^2u
 -\nabla\partial_t^2p.
\end{aligned}
\]
\begin{samepage}
The pattern is
\[
\partial_t^{\,n+1}u
=\nu\Delta\partial_t^{\,n}u
-\sum_{a=0}^{n}\binom na
  \bigl(\partial_t^{\,a}u\!\cdot\nabla\bigr)
  \partial_t^{\,n-a}u
-\nabla\partial_t^{\,n}p.
\]
\end{samepage}
\begin{samepage}
Set
\[
V_n:=\partial_t^{\,n}u,
\qquad
Q_n:=\partial_t^{\,n}p,
\]
Taking the divergence and using incompressibility gives the pressure at that
order:
\[
-\Delta Q_n
=\sum_{a=0}^{n}\binom na
  \partial_i\partial_j
  \bigl((V_a)_i(V_{n-a})_j\bigr).
\]
\end{samepage}
The line to keep in view is \(V_{n+1}=\nu\Delta V_n-\cdots\). Moving one step
up in time brings in two spatial derivatives of the response below it.
Transport and pressure stay attached, but viscosity supplies the link from
temporal order to spatial Sobolev order.

At the critical height, set
\[
\Lambda=(-\Delta)^{1/2},
\qquad
E_s(t):=\frac12\|\Lambda^su(t)\|_2^2,
\qquad
H(t):=E_{1/2}(t).
\]
If \(\mathcal P_s\) denotes the complete transport--pressure contribution at
height \(s\), the energy identity is
\[
E_s'=\mathcal P_s-2\nu E_{s+1}.
\]
Starting from \(H=E_{1/2}\) and repeating that identity gives
\[
H^{(n)}
=(-2\nu)^nE_{n+1/2}
 +\sum_{j=0}^{n-1}(-2\nu)^j
   \partial_t^{\,n-1-j}\mathcal P_{j+1/2}.
\]
The first term has reached spatial height \(n+\tfrac12\); the sum carries every
transport--pressure term generated on the way. This is the opposite of the
blow-up picture. Rising through that temporal sequence was supposed to reveal
less spatial regularity. The equation exposes higher Sobolev height instead. As
you do that, you get higher order Sobolev embeddings:
\[
s>k+\frac32
\quad\Longrightarrow\quad
H^s(\mathbb R^3)\hookrightarrow C^k(\mathbb R^3),
\qquad
\bigcap_{s<\infty}H^s(\mathbb R^3)
=H^\infty(\mathbb R^3)
\hookrightarrow C^\infty(\mathbb R^3).
\]

Seeing a higher energy in the identity is not enough. It has to remain finite
all the way to \(T_*\). Fix smooth divergence-free finite-energy data \(u_0\),
viscosity \(\nu>0\), and its maximal pressure-complete solution \((u,p)\) on
\([0,T_*)\). Assume \(T_*<\infty\). For finite \(N\) and \(M\), collect the
adjacent rows in
\[
\mathfrak R_{N,M}(u;T)
:=\sum_{n=0}^{N}\left(
  \|V_n\|_{L^2(0,T;H^{M+1})}^2
 +\|V_{n+1}\|_{L^2(0,T;H^{M-1})}^2
\right).
\]
Adding the corresponding \(Q_n\) and \(\nabla Q_n\) coordinates gives the
finite pressure-complete rectangle \(\mathcal R_{N,M}[u_0;T]\). The
manuscript's whole-terminal theorem bounds each adjacent pair on the full
interval:
\[
V_n\in L^2(0,T_*;H^{M+1}),
\qquad
V_{n+1}\in L^2(0,T_*;H^{M-1}),
\qquad 0\le n\le N.
\]
Restriction to \((0,T)\) can only decrease those norms, and there are only
finitely many rows, so
\[
\sup_{0<T<T_*}\mathfrak R_{N,M}(u;T)<\infty.
\]
\begin{samepage}
The bound may depend on \(N\) and \(M\). No uniform constant over all orders is
needed. At \(n=0\), for every finite \(m\), the adjacent pair gives
\[
u\in L^2(0,T_*;H^{m+1}),
\qquad
\partial_tu\in L^2(0,T_*;H^{m-1}),
\]
The Lions--Magenes trace theorem applied to
\(H^{m+1}\hookrightarrow H^m\hookrightarrow H^{m-1}\), gives a unique
trace at the endpoint:
\[
u\in C([0,T_*];H^m)
\qquad\text{for every finite }m.
\]
\end{samepage}

\begin{samepage}
Increasing \(N\) or \(M\) does not create a new solution. Every rectangle is a
finite part of the maximal history from \(u_0\), so overlapping rectangles
agree on every shared coordinate. Their projective limit is
\[
\mathcal I[u_0]
:=\varprojlim_{N,M<\infty}\mathcal R_{N,M}[u_0;T_*],
\qquad
u\in\bigcap_{r,m<\infty}H^r(0,T_*;H^m(\mathbb R^3)).
\]
\end{samepage}
The intersection is not one bound over infinitely many orders. Each coordinate
was bounded at a finite \((N,M)\); compatibility says that all of those finite
bounds belong to one velocity--pressure history.

The zeroth temporal row, \(\mathscr S_0:=\pi_0\mathcal I[u_0]\), now has one
terminal value at every finite spatial height:
\[
u_*:=u(T_*)
\in\bigcap_{m<\infty}H^m_\sigma(\mathbb R^3)
=H^\infty_\sigma(\mathbb R^3)
\hookrightarrow C^\infty(\mathbb R^3).
\]
The recurrence and pressure equation recover every finite coordinate of the
terminal jet from \(u_*\). The uniform \(H^3\) row gives a local existence time
\(\delta_0>0\) that works for \(u_*\) and for every sufficiently late slice
\(u(t_0)\). Choose \(t_0<T_*\) with \(T_*-t_0<\delta_0/2\), and restart from
\(u(t_0)\). Uniqueness makes this restart identical to the original solution
on \([t_0,T_*)\), so it reaches \(T_*\) with value \(u_*\). To the right of
\(T_*\), uniqueness identifies it with the solution restarted from \(u_*\).
The two restarts join into one smooth pressure-complete solution beyond
\(T_*\), contradicting maximality. Thus \(T_*=\infty\).

Once \(V_r\in C_tH_x^m\) is available for every finite \(r\) and \(m\), any
finite downstream estimate comes from choosing
\[
m>k+\frac32-\frac3p,
\qquad
2\le p\le\infty,
\]
which gives, on every finite time interval,
\[
V_r\in C_tH_x^m
\hookrightarrow L_t^qW_x^{k,p},
\qquad 1\le q\le\infty.
\]

\[
u_0
\longrightarrow \{(V_n,Q_n)\}_{n\ge0}
\longrightarrow \{\mathcal R_{N,M}[u_0]\}_{N,M<\infty}
\longrightarrow \mathcal I[u_0]
\longrightarrow \mathscr S_0
\longrightarrow u_*\in H^\infty
\longrightarrow T_*=\infty.
\]
